\section{Incidence and period descent without global labels}\label{sec:descent}
\subsection{Complete images and their incidence}
For a convex body $C$ let $p_C$ be its support function after translation
to its Steiner point. Congruence includes orthogonal transformations of
either determinant. A strictly convex real-analytic boundary has a
real-analytic support function on the circle: positive curvature makes
the normal-angle map analytically invertible. Two such bodies sharing
an open boundary arc in one frame have equal support functions on an
open normal-angle interval, hence everywhere by the identity theorem.
Thus Theorem~\ref{thm:local-main} determines each entire analytic pair
in its relative placement. No smooth-jet identity is used for this step.

\begin{lemma}[Intrinsic incidence recovery]\label{lem:incidence}
For an asymmetric, shape-separated table, a connected covering unlabelled
collection determines its incidence multigraph and places one representative
body of each vertex along any chosen spanning tree, up to a common isometry.
Every non-tree edge determines a unique translation $d_e$ between its
target image and the chosen representative. Each $d_e$ belongs to the
unknown period lattice. The construction uses no integer edge labels.
\end{lemma}
\begin{proof}
Reconstruct the two complete images from each local record. Partition
all image occurrences by Euclidean congruence. Shape separation means
that two occurrences are in the same class precisely when their physical
bodies are translates of the same obstacle representative. Coverage makes
the classes exactly the vertices; the two image slots of each record
give its incident vertices, including loops and repeated edges.

If $C$ and $C'$ are congruent asymmetric bodies, there is exactly one
isometry carrying $C$ to $C'$. Indeed two such maps differ by a symmetry
of $C$. Place one root image arbitrarily. Along an outward tree edge,
match its source image to the already placed source representative and
apply the same isometry to the opposite image. This places the child.
For backward traversal interchange the two images; no additional
observation is required. A change of the local reversal representative
is cancelled by composing its unique matching map with that reversal.

On a non-tree edge match the source in the same way. Let $\widehat C_w$
be the target image and $C_w$ the chosen target representative. In a
physical table they differ by a period translation. This translation is
unique since a nonempty compact set has no nonzero translational symmetry.
It is computed from Steiner points:
\begin{equation}\label{eq:cycle-displacement}
 d_e=s(\widehat C_w)-s(C_w),\qquad \widehat C_w=C_w+d_e.
\end{equation}
The second equality is a consistency condition, not an extra matching
choice. Tree placement chooses physical copies of the true obstacles;
therefore $d_e\in\Lambda$. All output changes under root placement by
one common orthogonal transformation and translation only.
\end{proof}

The graph is recovered from shapes, not from a comparison of curvature
values at isolated contacts. Pairwise noncongruence is essential for
this particular incidence argument. It is not an assumption that a
shape catalogue has been supplied to the observer. On the exact class,
congruence is an equality test on recovered analytic functions; no
polynomial-time test of such equality is asserted.

\subsection{The normalization determines the covolume}
Every record supplies the same $A$ in \eqref{eq:ratio-main}. Since the
incidence classes contain exactly one obstacle representative per period
cell, the recovered images give
\[
             V=A+\sum_i\area(C_i)=\covol\Lambda.
\]
This identity would not hold with a sum over image occurrences, which
would count repeated observations several times. It is also why the
shape-separated hypothesis cannot simply be discarded while removing
obstacle labels. Areas can be computed intrinsically from the images;
for example
\begin{equation}\label{eq:support-area}
 \area(C)=\frac12\int_0^{2\pi}(p_C^2-(p_C')^2)\dd\theta.
\end{equation}
The formula is unchanged by translating $C$.

\begin{lemma}[Finite-index lattice classification]\label{lem:overlattice}
Let $D$ be an invertible real $2\times2$ matrix and $V>0$. The lattices
$\Lambda$ containing $D\Z^2$ and having covolume $V$ form the list
\eqref{eq:hnf-main}, provided $n=|\det D|/V$ is a positive integer.
Otherwise there is no such lattice. The list has exactly $\sigma_1(n)$
distinct lattices. An additional vector $d_e$ belongs to $\Lambda_B$
if and only if $BD^{-1}d_e\in\Z^2$.
\end{lemma}
\begin{proof}
The subgroup $D\Z^2\subset\Lambda$ has index $n$, by the ratio of
fundamental parallelogram areas. Normalize by $D^{-1}$ and write
$L=D^{-1}\Lambda\supset\Z^2$. Its dual
$L^*=\{z:z\cdot L\subset\Z\}$ is a sublattice of $\Z^2$ of
index $n$. Conversely every such sublattice is the dual of a unique $L$.

For completeness, a sublattice $H\subset\Z^2$ of finite index has
a unique row basis $(a,b),(0,d)$ with $a,d>0$ and $0\le b<d$.
Its projection on the first coordinate is $a\Z$, its vertical intersection
is $\{0\}\times d\Z$, and subtraction of vertical multiples gives the
unique representative $b$. Every element then subtracts to zero using
these two rows. Its index is $ad$. Thus $H$ is the row span of the
matrix $B$ in \eqref{eq:hnf-main}, and its dual is $B^{-1}\Z^2$.
For each $d\mid n$ there are $d$ choices of $b$. This proves the count,
uniqueness of the list, and the stated membership test.
\end{proof}

\begin{proof}[Proof of Theorem~\ref{thm:descent-main}]
Lemma~\ref{lem:incidence} recovers all representative images and cycle
displacements. The normalization recovers $V$. Two independent $d_e$
generate a sublattice $D\Z^2$ of the true period lattice, so
Lemma~\ref{lem:overlattice} applies. For each listed lattice use the
already placed representative bodies. Test all remaining cycle vectors,
disjointness of translates, area and the selected physical channels.
A true table producing the records appears on the list and passes every
test. There is no additional continuous placement parameter. The list
therefore bounds the number of compatible tables. If $n=1$ it consists
only of $D\Z^2$.

The outcome is independent of the initial local reflections and root
placement by Lemma~\ref{lem:incidence}. A different spanning tree changes
the chosen copies and the cycle generators but not any physical table
passing the tests. One can enumerate the finitely many tree and cycle-pair
choices to find a determinant equal to $V$; the observer need not be told
which pair has this property. This proves the unique case as well.
\end{proof}

\begin{remark}[The unmarked period lattice]
No preferred basis of $\Lambda$ is recovered or required. In the primitive
case $D$ is one recovered basis, up to the usual unimodular change. On the
shape-separated class $\Lambda$ is in fact the full translational period
group: a translation preserving the table must preserve each shape class,
and hence sends $C_i$ to $C_i+\lambda$ for some $\lambda\in\Lambda$.
Compactness forces that translation to be $\lambda$. Thus the conclusion
does not determine an artificially chosen nonminimal covering torus.
\end{remark}

\subsection{Physical primitive classes and finite-index ambiguity}
\begin{proposition}[Nonempty open primitive classes]\label{prop:primitive}
For every positive integer $r$ there are nonempty open analytic families
with $r$ pairwise noncongruent asymmetric obstacle representatives and
a connected covering collection satisfying $|\det D|=V$.
\end{proposition}
\begin{proof}
Use small strictly convex bodies with centered support functions
\[
 p_i(\theta)=r_i\{1+\varepsilon\cos(2\theta)
                    +\varepsilon\cos(3\theta+\pi/4)\},
                \qquad 0<|\varepsilon|<1/11,
\]
where the positive $r_i$ are distinct. The curvature radius $p_i+p_i''$
is positive. A rotation preserving both Fourier modes has angle zero
modulo $2\pi$. A reflection about angle $a$ would require both
$2a\in\pi\Z$ and $3a+\pi/4\in\pi\Z$, which is impossible.
Distinct scales are noncongruent.

Take a rectangular lattice with periods $(L,0),(0,M)$. Place one small
body at the origin and the others at generic points strictly inside the
cell. Choose a star joining the root to these points, and include the
root's horizontal and vertical neighboring-copy channels. Generic points
ensure that no other center lies on any chosen segment, including relevant
nearby translated centers. Since there are only finitely many such
segments and nearby centers, sufficiently small bodies leave positive
clearance. The normal channels persist by the nondegenerate stationary
minimum. The star is a spanning tree, and the two root loops yield exactly
$(L,0)$ and $(0,M)$, whose determinant equals the period covolume.
Positive clearance, convexity, asymmetry and pairwise noncongruence persist
in a sufficiently small analytic neighborhood. The loop labels need not
be observed: their geometric displacements remain the two periods.
\end{proof}

\begin{proposition}[Physical finite-index ambiguity]\label{prop:ambiguity}
For each odd prime $p$ there are $p-1$ nonisometric analytic periodic
tables with identical unlabelled two-density records, identical free area,
a connected covering graph and two independent reconstructed displacements,
for which the calibrated index is $p$.
\end{proposition}
\begin{proof}
Fix one asymmetric analytic strictly convex body $C$ contained in a disk
of radius $r$, and choose $L,M>8pr$. Put $d_1=(L,0)$, $d_2=(0,M)$
and, for $j=1,\ldots,p-1$, set
\[
 \Lambda_j=\Z d_1+\Z d_2+
                      \Z\frac{d_1+j d_2}{p}.
\]
These are distinct superlattices of $\Z d_1+\Z d_2$ of index $p$;
each has covolume $LM/p$. Their nonzero vectors have length at least
$\min(L,M)/p$, so the copies of $C$ are disjoint. All tables have free
area $A=LM/p-\area(C)>0$.

Observe only the channels from $C$ to $C+d_1$ and from $C$ to $C+d_2$.
A center in $\Lambda_j$ on the horizontal axis is an integral multiple
of $d_1$, because $j$ is invertible modulo $p$. A center on the vertical
axis is an integral multiple of $d_2$. Off the horizontal axis its
vertical distance is at least $M/p$, and off the vertical axis its
horizontal distance is at least $L/p$. Thus the two selected pair
corridors, lying within radius $r$ of their axis segments, are separated
from every other body by a positive common margin. Their normal-channel
geometries are independent of $j$. On sufficiently small common patches
and time collars, Lemma~\ref{lem:density} gives exactly the same densities:
both $W$ and $A$ agree. This is equality of physical laws, not merely
of finite jets. The recovered graph has one vertex and two loops, and
$|\det[d_1\ d_2]|/V=p$.

Finally an isometry between two tables sends a connected obstacle component
to a translate of $C$. Its orthogonal part must preserve the centered body,
and is therefore the identity. A translation conjugates neither period
group to a different one. Since $\Lambda_j$ are distinct and are the full
period groups, the tables are nonisometric.
\end{proof}

The proposition locates the finite obstruction rather than replacing the
rigidity theorem by nonidentifiability. The primitive theorem gives unique
reconstruction on an open physical class; the general theorem exhausts
all remaining period choices by a finite arithmetic list.
